\chapter{Protein in Urine (Proteinuria)}

\section*{Definition}
Proteinuria is an abnormally increased amount of protein in urine. It may be temporary or indicate kidney disease, diabetes, hypertension, infection, or pregnancy-related disease.

\section*{When to See a Doctor}
Seek urgent evaluation for Pregnancy with high blood pressure, headache, visual symptoms, swelling, reduced urine, breathlessness, or rapidly worsening kidney function. Arrange routine or prompt assessment for persistent, recurrent, unexplained, or function-limiting symptoms.

\section*{How It Develops}
The symptom reflects irritation, inflammation, obstruction, altered secretion, infection, or dysfunction in the relevant organ system. The pattern, duration, associated symptoms, exposures, and examination findings determine whether it is self-limited or requires investigation.

\section*{Causes}
\begin{itemize}
  \item Diabetes, hypertension, glomerular disease, fever, strenuous exercise, dehydration, infection, preeclampsia, and medications.
  \item Medication effects, environmental exposures, and chronic disease may contribute.
  \item Less common but important causes include malignancy, vascular disease, or organ failure when symptoms persist or progress.
\end{itemize}

\section*{Related Symptoms}
\begin{itemize}
  \item Pain, fever, fatigue, nausea, or changes in normal organ function
  \item Bleeding, weight change, shortness of breath, weakness, or neurologic symptoms when a serious cause is present
\end{itemize}

\section*{Medical Evaluation}
\begin{itemize}
  \item History addresses onset, duration, severity, triggers, exposures, medications, medical conditions, pregnancy status when relevant, and warning symptoms.
  \item Examination focuses on vital signs and the organ systems suggested by the symptom.
  \item Testing may include blood or urine studies, cultures, imaging, physiologic testing, or specialist examination according to the clinical pattern.
\end{itemize}

\section*{Treatment Options}
Treatment is directed at the cause and may include supportive care, targeted medication, treatment of infection or inflammation, correction of metabolic disease, or specialist procedures. Persistent symptoms require reassessment rather than indefinite symptomatic treatment.

\section*{Self-Care and Home Management}
Avoid known triggers, maintain reasonable hydration, record timing and associated symptoms, and use only treatments appropriate for the suspected cause. Do not use leftover antibiotics or delay urgent care because symptoms temporarily improve.

\section*{References}
\begin{thebibliography}{99}
\bibitem{protein_in_urine:kdigo} KDIGO 2021 Clinical Practice Guideline for the Management of Glomerular Diseases. \textit{Kidney Int}. 2021;100(4S):S1--S276.
\bibitem{protein_in_urine:nice-ckd} National Institute for Health and Care Excellence. Chronic kidney disease: assessment and management. NICE guideline NG203; 2021 (updated 2025).
\bibitem{protein_in_urine:harrison} Jameson JL, et al. \textit{Harrison's Principles of Internal Medicine}. 21st ed. McGraw-Hill; 2022.
\end{thebibliography}
